\section{Full problem setup}
\label{app:setup}

We use the size convention of Section~\ref{sec:main-setup}: $N$ is fixed
within each finite statement, and size dependence is suppressed, including
in $\TargetLaw,\VocabSize,\EdgeSignal$ and the accuracy parameters.
In rate statements, parameters declared constant are independent of $N$;
every public sequence and its deterministic remainder is fixed uniformly
over the class before any hidden target or oracle is chosen.

\paragraph{Reading route to the main claims.}
The formal setup below fixes the objects used by all proofs.
The following routes give the dependencies of the main claims:
\begin{itemize}
 \item \emph{Lower bound in \HellingerCase, \cref{main:thm:lower-envelope}.}
 Appendix~\ref{app:results} defines the matching family and states its
 finite lower bound.  Appendix~\ref{app:lower-proof} proves the bound and
 verifies class inclusion in \cref{app:prop:hard-family-inclusion}.
 Then the adjacent
 \cref{app:cor:signal-calibrated-finite-lower,app:cor:fixed-accuracy-main}
 in Appendix~\ref{app:benchmark-specializations} give the finite
 calibration and the fixed-accuracy deduction of the main theorem.
 \item \emph{Upper bounds in both cases, \cref{main:thm:upper-curve}.}
 Appendix~\ref{app:algorithm} specifies the executable sampler.
 Appendix~\ref{app:upper-proof} proves the shared screen, peeling, and
 resource bounds, and the TV guarantee in \HellingerCase\
 (\cref{app:thm:finite-upper}).
 Appendix~\ref{app:kl-extension} proves the forward-KL guarantee in
 \KLCase\ (\cref{app:thm:finite-kl-upper}), using
 \cref{app:lem:adaptive-joint-kl}.
 \Cref{app:lem:benchmark-feasibility,app:cor:fixed-accuracy-main}
 verify the finite calibration and derive \cref{main:thm:upper-curve}
 with accuracy fixed before $N\to\infty$.
 \ifincludeexponents
 \Cref{app:cor:benchmark-hellinger,app:cor:benchmark-kl} give the further
 joint-limit consequences when accuracy also varies with $N$.
 \fi
 \item \emph{Comparison in Table~\ref{tab:resource-comparison}.}
 Appendix~\ref{app:comparison} proves the row-to-state accounting and
 the path's TC/DTC scales; Appendix~\ref{app:oracle-comparison} records
 the assumptions of the source results being substituted.
 \item \emph{Finite-size experiments, Section~\ref{sec:main-experiments}.}
 Appendix~\ref{app:experiments} gives the targets, accuracy test, tuning,
 resource accounting, and evaluation protocol.
\end{itemize}
The elementary auxiliary facts in Appendix~\ref{app:auxiliary} are
referenced at their point of use.  That appendix also collects the
supplementary selector-invariance result.
\ifincludeexponents
Appendix~\ref{app:independent-signal} gives the additional rates when
the signal floor is specified independently; these are not needed for
the two fixed-accuracy main theorems.
\fi

Short intuition paragraphs precede the main proof steps.
There, $\simeq$ compares scales up to fixed positive multiplicative
constants in the stated regime.  An upper bound alone is written with
$\lesssim$ or an explicit inequality, and subpolynomial factors are
displayed when relevant.  All formal statements retain their finite
constants and domains.

\begin{proposalcontext}{Supplementary context}
Unnumbered remarks give supplementary context beyond the proof route above.
Assumptions, resource definitions, conditioning arguments, and restrictions
needed to apply a result are stated in the main exposition.
\end{proposalcontext}

For each $N\ge 3$, let $\PositionSet:=\{1,\ldots,N\}$ be the set of positions and
let $\Vocab$ be a finite vocabulary of size $\VocabSize\ge 2$, equipped with
a public total order $\TokenOrder$.  For a finite set $\mathcal A$, write
$\Delta(\mathcal A)$ for its probability simplex and
$\Delta^\circ(\mathcal A)$ for the strictly positive part.  For
\(p,q\in\Delta(\mathcal A)\), we use the normalizations
\begin{equation}
 \dTV(p,q)
 :=\frac12\sum_{a\in\mathcal A}|p(a)-q(a)|,
 \qquad
 \sqHellinger(p,q)
 :=1-\sum_{a\in\mathcal A}\sqrt{p(a)q(a)}.
 \label{app:eq:distances}
\end{equation}
For \(q\in\Delta^\circ(\mathcal A)\), also write
\begin{equation}
 \ChiSquareDivergence{p}{q}
 :=\sum_{a\in\mathcal A}\frac{(p(a)-q(a))^2}{q(a)}.
 \label{app:eq:chi-square-divergence}
\end{equation}
Also define $\KLDivergence{p}{q}:=\sum_a p(a)\log(p(a)/q(a))$, with
natural logarithms and the convention $0\log(0/q(a))=0$.
For \(x\in\mathbb R\), let \(\PositivePart{x}:=\max\{x,0\}\).

All quantities explicitly declared public below, together with the algorithm
and the law of its internal random seed, are fixed before the target instance
is chosen.  The hidden target law and, subsequently, the frozen oracle may be
chosen adversarially.  Local vocabulary banks $\TokenBank{i}$, tail
representatives, drafts, and graph certificates are constructed from oracle
replies by the algorithm.

\subsection{Target law and hidden forest}
\label{app:setup-target}

We expand \cref{main:def:forest-target,main:ass:forest-structure}.

\begin{definition}[Forest-structured target law]
\label{app:def:target-law}
Let $X=(X_1,\ldots,X_N)\in\Vocab^N$ have a hidden strictly positive law
$\TargetLaw\in\Delta^\circ(\Vocab^N)$, and let
$\TargetForest=(\PositionSet,\TargetEdges)$ be a hidden undirected forest.  The exact position
marginal is
\begin{equation}
 \PosMarginal{i}:=\mathcal L_{\TargetLaw}(X_i)
 \in\Delta^\circ(\Vocab),
 \qquad i\in\PositionSet.
 \label{app:eq:position-marginal}
\end{equation}
For each orientation $i\to j$ of an edge $\{i,j\}\in\TargetEdges$, let
$\EndpointKernel{i}{j}(\cdot\mid a)\in\Delta^\circ(\Vocab)$ be an endpoint
kernel for every $a\in\Vocab$.
\end{definition}

\begin{assumption}[Structural forest law (S0--C0)]
\label{app:ass:forest-structure}
\textbf{Forest factorization (S0).}
Fix one hidden reference root in each connected component.  If
$\TargetRoots$ is the resulting root set and $\ParentOf{j}$ is the
induced parent of a nonroot vertex $j$, assume
\begin{equation}
 \TargetLaw(x)
 =\prod_{u\in\TargetRoots}\PosMarginal{u}(x_u)
  \prod_{j\notin\TargetRoots}
  \EndpointKernel{\ParentOf{j}}{j}
  (x_j\mid x_{\ParentOf{j}}).
 \label{app:eq:S0}
\end{equation}

\textbf{Endpoint compatibility (C0).}
For every $\{i,j\}\in\TargetEdges$ and $a,b\in\Vocab$, assume
\begin{equation}
 \PosMarginal{i}(a)\EndpointKernel{i}{j}(b\mid a)
 =\PosMarginal{j}(b)\EndpointKernel{j}{i}(a\mid b).
 \label{app:eq:C0}
\end{equation}
Moving a root across one edge replaces the left-hand side of
\ref{app:eq:C0} by its right-hand side and leaves every other factor
unchanged.  Iterating this replacement along the unique path between two roots
shows that (S0)--(C0) yield the same factorization for every componentwise root
choice.
\end{assumption}

To obtain these kernels from \cref{main:ass:forest-structure}, set
$\EndpointKernel{i}{j}(b\mid a):=\TargetLaw(X_j=b\mid X_i=a)$
for each oriented edge. Then (S0) is the main-text factorization, and
(C0) follows from
\[
 \PosMarginal{i}(a)\EndpointKernel{i}{j}(b\mid a)
 =\TargetLaw(X_i=a,X_j=b)
 =\PosMarginal{j}(b)\EndpointKernel{j}{i}(a\mid b).
\]

For a committed set $G\subseteq\PositionSet$ and realized values
$x_G\in\Vocab^G$, define the history, uncommitted set, and residual
forest by
\begin{equation}
 \History:=(G,x_G),
 \qquad
 \ResidualPositions:=\PositionSet\setminus G,
 \qquad
 \ResidualForest:=\TargetForest[\ResidualPositions].
 \label{app:eq:history}
\end{equation}
Strict positivity makes every such history admissible.
For any graph $H=(W,E_H)$, we write
$\mathcal N_H(v):=\{u\in W:\{u,v\}\in E_H\}$ for its neighborhood and
$d_H(v):=|\mathcal N_H(v)|$ for its degree.  The induced graph $H[S]$,
for $S\subseteq W$, has vertex set $S$ and edge set
$\{\{u,v\}\in E_H:u,v\in S\}$; $V(T)$ denotes the vertex set of a
connected component $T$.

\subsection{Response regularity}
\label{app:setup-response}

\textcolor{red}{The response quantities below correspond to
\cref{main:def:response-quantities}; the general regularity assumptions
correspond to the RT, UEN, and RF items of
\cref{main:ass:regularity-scaling}.}
For simplicity, the main text and \cref{app:cor:fixed-accuracy-main}
use $\ResponseExponent=\ResponseConstant=\FrequencyConstant=1$.

\begin{definition}[Marginal tails and intrinsic responses]
\label{app:def:response-quantities}
For a source position $i\in\PositionSet$ and threshold $t\in[0,1]$, define the exact
marginal tail
\begin{equation}
 \TailSet{i}{t}
 :=\{a\in\Vocab:\PosMarginal{i}(a)\le t\}.
 \label{app:eq:tail}
\end{equation}
For distinct positions $i,j$ and a complete boundary
$z\in\Vocab^{\PositionSet\setminus\{i,j\}}$, define the boundary-conditioned
row
\begin{equation}
 \BoundaryRow{i}{j}{a}{z}
 :=\mathcal L_{\TargetLaw}\!\left(
 X_j\mid X_i=a,
 X_{\PositionSet\setminus\{i,j\}}=z
 \right).
 \label{app:eq:boundary-row}
\end{equation}
For an edge $e=\{u,v\}\in\TargetEdges$ and such a complete boundary, define
the directed intrinsic response
\begin{equation}
 \DirectedResponse{v}{u}{z}
 :=\max_{a,b\in\Vocab}
 \dTV\bigl(
 \BoundaryRow{v}{u}{a}{z},
 \BoundaryRow{v}{u}{b}{z}
 \bigr).
 \label{app:eq:intrinsic-response}
\end{equation}
\end{definition}

\begin{assumption}[Response regularity (RT--UEN)]
\label{app:ass:response-regularity}
Let $\ResponseExponent>0$, $\ResponseConstant\in[0,\infty)$, and
$\EdgeSignal\in(0,1]$ be public.

\textbf{Power tail-response continuity (RT).}
For all distinct $i,u\in\PositionSet$, every $t\in[0,1]$, and every complete assignment
$z\in\Vocab^{\PositionSet\setminus\{i,u\}}$, assume
\begin{equation}
 \max_{a,b\in\TailSet{i}{t}}
 \dTV\bigl(
 \BoundaryRow{i}{u}{a}{z},
 \BoundaryRow{i}{u}{b}{z}
 \bigr)
 \le \ResponseConstant t^{\ResponseExponent}.
 \label{app:eq:RT}
\end{equation}
The maximum is defined to be zero when the tail is empty or a singleton.

\textbf{Two-sided edge nondegeneracy (UEN).}
For every edge
$e=\{u,v\}\in\TargetEdges$ and every complete boundary $z$, assume
\begin{equation}
 \min\bigl\{
  \DirectedResponse{v}{u}{z},
  \DirectedResponse{u}{v}{z}
 \bigr\}
 \ge \EdgeSignal.
 \label{app:eq:UEN}
\end{equation}
Only the scale $\EdgeSignal$ is public; the forest, edge direction,
boundary, and maximizing token pair remain hidden.
\end{assumption}

\begin{assumption}[Asymptotic vocabulary and frequency envelope (RF)]
\label{app:ass:rate-regime}
For asymptotic rate statements, let $\VocabGrowthExponent>0$,
$\RankTailExponent>1$, and $\FrequencyConstant\ge1$ be public constants, and
assume $\VocabSize=N^{\VocabGrowthExponent+o(1)}$.  Order the vocabulary at each position by exact
marginal mass,
\begin{equation}
 \PosMarginal{i}(\RankedToken{i}{1})\ge\cdots\ge
 \PosMarginal{i}(\RankedToken{i}{\VocabSize}),
 \label{app:eq:hidden-ranks}
\end{equation}
with ties broken by $\TokenOrder$, and assume uniformly in $i\in\PositionSet$ and
$k\in[\VocabSize]$ that
\begin{equation}
 \PosMarginal{i}(\RankedToken{i}{k})
 \le \FrequencyConstant
      (k^{-\RankTailExponent}+\VocabSize^{-1}).
 \label{app:eq:RF}
\end{equation}
No lower Zipf bound or adjacent-rank separation is assumed.  This assumption
is used only to turn finite bounds into uniform asymptotic rates; it is not
needed for finite correctness once threshold feasibility and the realized bank
sizes are fixed.
\end{assumption}

% Please keep the two case descriptions together with their introductory heading.
\begin{samepage}
\subsection{Frozen conditional oracle}
\label{app:setup-oracle}

The two cases of Assumption~\ref{main:ass:oracle-accuracy} share the frozen oracle
interface below, with different accuracy conditions and output objectives:
\begin{itemize}
 \item \emph{\HellingerCaseTitle\ (Hellinger/TV).}
 The uniform squared-Hellinger condition is (A2) below; the output objective
 is the seed-averaged TV risk in \textcolor{red}{\cref{app:def:output-risk}}.
 The normalized Hellinger radius $\OracleRadius$ remains independent
 of the output-TV tolerance $\TargetTolerance$.
 \item \emph{\KLCaseTitle\ (forward KL).}
 \Cref{app:def:kl-oracle-risk} defines the uniform forward row-KL
 condition and the seed-averaged forward-KL output objective.
\end{itemize}
\end{samepage}
The main text uses $\OracleRadius=\TargetTolerance/2$ in \HellingerCase\
and $\KLTargetTolerance=\TargetTolerance^2/2$ in \KLCase, with the latter
calibrated by \cref{app:cor:kl-epsilon}.
\Cref{app:cor:fixed-accuracy-main} gives both main-text specializations.

\begin{definition}[Frozen conditional oracle]
\label{app:def:frozen-oracle}
As in \cref{main:def:frozen-oracle}, let $\MASK\notin\Vocab$. For a masked state
$y\in(\Vocab\cup\{\MASK\})^N$, define
\begin{equation}
 \MaskSet{y}:=\{i:y_i=\MASK\},
 \qquad
 \ObservedSet{y}:=\PositionSet\setminus\MaskSet{y}.
 \label{app:eq:masked-revealed}
\end{equation}
For $j\in\MaskSet{y}$, the hidden exact singleton conditional is
\begin{equation}
 \ExactRowMass{y}{j}{a}
 :=\TargetLaw\bigl(X_j=a\mid
 X_{\ObservedSet{y}}=y_{\ObservedSet{y}}\bigr),
 \qquad a\in\Vocab.
 \label{app:eq:exact-row}
\end{equation}
The all-mask state is $\AllMaskState:=(\MASK,\ldots,\MASK)$, so
$\ExactRow{\AllMaskState}{i}=\PosMarginal{i}$.

The oracle is one deterministic frozen map
\begin{equation}
 \FrozenOracle:\{(y,j):y\in(\Vocab\cup\{\MASK\})^N,
                 \ j\in\MaskSet{y}\}
 \longrightarrow\Delta^\circ(\Vocab).
 \label{app:eq:oracle-type}
\end{equation}
Equivalently, write
$\FrozenOracle=\{\FrozenOracle_j\}_{j\in\PositionSet}$, with
$\FrozenOracle(y,j)=\OracleRow{y}{j}$ and
$\OracleRowMass{y}{j}{a}$ its probability mass at token $a$.

\emph{\textcolor{red}{One submission, a family of probability vectors.}}
One oracle evaluation submits a single masked state $y$ and returns full
vocabulary probability vectors at the requested masked positions:
\[
 (y,S)\ \xrightarrow{\text{one masked-state submission}}\quad
 \bigl\{\OracleRow{y}{j}\bigr\}_{j\in S},
 \qquad S\subseteq\MaskSet{y}.
\]
Here $S$ selects the readouts; each returned row is the whole vector
$\OracleRow{y}{j}=(\OracleRowMass{y}{j}{a})_{a\in\Vocab}$.
The charge is one submission, regardless of $|S|$ or how many tests use
the returned probabilities. Every occurrence of the same $(y,j)$ returns
the same row, across readout sets and operations; repeated submissions
still count as separate evaluations.
Probes and commits in \cref{app:def:admissible-algorithm} use this same
interface and the same frozen rows.
\end{definition}

\begin{assumption}[\HellingerCaseTitle: frozen-oracle accuracy (A2)]
\label{app:ass:oracle-accuracy}
For a public $\OracleRadius\in[0,1]$, assume uniformly over all valid $(y,j)$ that
\begin{equation}
 \sqHellinger\bigl(\ExactRow{y}{j},\OracleRow{y}{j}\bigr)
 \le \frac{(\OracleRadius)^2}{2N}.
 \label{app:eq:A2}
\end{equation}
Oracle errors may be adversarially correlated across states, positions, and
calls; no independence, unbiasedness, concentration, or fresh-noise condition
is imposed.  Let $\OracleClass(\TargetLaw;\OracleRadius)$ denote the class
of frozen maps satisfying (A2).
\end{assumption}

\begin{remark}[Oracle metric calibration]
\label{app:rem:oracle-metric-calibration}
The parameter $\OracleRadius$ is a normalized oracle budget: the per-row
Hellinger distance in (A2) is at most $\OracleRadius/\sqrt{2N}$, not
$\OracleRadius$.  Lemma~\ref{app:lem:finite-hellinger-calculus} therefore
gives the row-TV bound $\OracleRadius/\sqrt N$ used by the screen; this is
a consequence of (A2), not an additional assumption.  Conversely, that
row-TV bound alone does not imply (A2).  For example, for sufficiently small
$e>0$, the positive binary laws $(1-2e,2e)$ and $(1-e,e)$ have TV distance
$e$ but squared Hellinger distance at least
$e(\sqrt2-1)^2/2>e^2/2$.

By Lemma~\ref{app:lem:kl-hellinger}, either of the following uniform bounds
is sufficient for (A2):
\begin{equation}
 \sup_{(y,j)}\KLDivergence{\ExactRow{y}{j}}{\OracleRow{y}{j}}
 \le\frac{(\OracleRadius)^2}{N}
 \quad\text{or}\quad
 \sup_{(y,j)}\KLDivergence{\OracleRow{y}{j}}{\ExactRow{y}{j}}
 \le\frac{(\OracleRadius)^2}{N}.
 \label{app:eq:row-kl-sufficient}
\end{equation}
These are \emph{sufficient conditions}, not necessary bounds on the KL error of an
oracle satisfying (A2): Remark~\ref{app:rem:no-hellinger-kl-converse} gives
a strictly positive counterexample to the converse.  Imposing either KL
condition leaves the upper theorem valid with its stated output-TV bound;
it does not assert an output-KL bound for the randomized sampler.
\end{remark}

\subsection{Algorithms, risk, and resources}
\label{app:setup-algorithms}

Both operations below use the row family of
\cref{app:def:frozen-oracle}. A \emph{probe} inspects its probabilities
without changing the committed history; a \emph{commit} samples a chosen
batch from these rows and permanently fixes the sampled values.
The algorithm requirements and resource counts below apply to both cases.
They give the full interface of \cref{main:def:algorithm-risk-resources}.

\begin{definition}[\textcolor{red}{Admissible algorithms}]
\label{app:def:admissible-algorithm}
\textbf{Algorithm requirements.}
An admissible algorithm accesses
$\FrozenOracle=\{\FrozenOracle_j\}_{j\in\PositionSet}$ adaptively and
satisfies the following conditions.
\begin{itemize}
\item \emph{Observable decisions and independent randomness.}
Let $\ControllerSeed$ be decision-rule randomness independent of the hidden
instance and of the primitive random sources used for product-commit draws.
The \emph{decision rule} (also called the controller) uses public data,
previous oracle replies, committed values, and its internal state and seed
to choose the next operation (probe or commit), its submitted state, and
any commit batch. Once the seed is fixed,
these choices are deterministic functions of the observable history.
The sampled token values may depend on $\ControllerSeed$ through the
decision rule's choices, but the primitive commit-draw randomness is
separate from that seed. The forest, exact marginals, endpoint kernels,
hidden ranks, edge directions, and response witnesses remain hidden.

\item \emph{No remasking and history-compatible submissions.}
At a realized history $\History=(G,x_G)$, every submitted state must agree
with the committed values: $y_i=x_i$ for all $i\in G$.
Committed positions stay revealed in all subsequent submissions, with the
same values. Uncommitted positions in $\ResidualPositions$ may be masked
or temporarily assigned vocabulary tokens for a probe.

\item \emph{Probes (counterfactual submissions).}
\textcolor{red}{Choose a history-compatible state $y$ and a subset of masked
positions $S\subseteq\MaskSet{y}$. One submission $(y,S)$ returns
$\{\OracleRow{y}{j}\}_{j\in S}$ together, allowing the algorithm to inspect
all probabilities $\{\OracleRowMass{y}{j}{a}:j\in S,\ a\in\Vocab\}$.}
This operation updates the observable transcript but leaves
$\History=(G,x_G)$ unchanged; temporary assignments are not committed.
Probes may be adaptive and interleaved with commits.
They may also use $\ObservedSet{y}=G$, with every residual position masked:
additional hypothetical reveals are optional.

\item \emph{Irreversible product commits.}
Let $t$ index commits only, and let $G_{t-1}$ be the positions committed
before round $t$, with $G_0=\emptyset$.
Choose a nonempty batch $B_t\subseteq\ResidualPositionsAt{G_{t-1}}$
from the observable past and $\ControllerSeed$, \emph{before} receiving
this commit's oracle reply. Submit the state
\begin{equation}
 (y_t)_i=
 \begin{cases}
  x_i, & i\in G_{t-1},\\
  \MASK, & i\in\ResidualPositionsAt{G_{t-1}}.
 \end{cases}
 \label{app:eq:commit-state}
\end{equation}
\textcolor{red}{One submission $(y_t,B_t)$ returns the probability vectors
$\{\OracleRow{y_t}{j}\}_{j\in B_t}$ from the same interface.}
Draw $z_{B_t}\in\Vocab^{B_t}$ according to
\begin{equation}
 \bigotimes_{j\in B_t}\OracleRow{y_t}{j},
 \label{app:eq:product-commit}
\end{equation}
and permanently update $G_t:=G_{t-1}\cup B_t$ and $x_{B_t}:=z_{B_t}$,
retaining the earlier values $x_{G_{t-1}}$.
The product in \eqref{app:eq:product-commit} draws the batch coordinates
independently conditional on the current history and selected batch.
Temporary probe assignments cannot replace this draw, and the draw is
not filtered by an accept--reject or verify-and-accept step.

\item \emph{Completion.}
The batches are disjoint and satisfy
$\bigcup_{t=1}^{\CommitRounds}B_t=\PositionSet$.
\end{itemize}

\smallskip
\noindent\textbf{Operation-indexed actions and visible history.}
\textcolor{red}{Index the submissions actually performed, probes and commits
together, by consecutive positive integers $s=1,2,\ldots$; repeated or
parallel submissions have separate indices. This operation index is
distinct from the rank-frequency exponent in \eqref{app:eq:RF}.}
Let $s_t$ be the operation index of commit $t$, and write $t_s=t$ when
$s=s_t$. Fix the decision-rule seed $w$.
The complete visible history immediately before operation $s$ is
$\OpVisibleHistory{s}$, initialized by the public data and seed-fixed
initial state, $\OpVisibleHistory{1}:=h_{\mathrm{init}}(w)$.
The selected action has the form
\begin{equation}
 A_s=\mathsf{Next}_w(\OpVisibleHistory{s})=
 \begin{cases}
  (\mathsf{query},\OpSubmittedState{s},J_s),
    &\text{probe},\\
  (\mathsf{commit},\OpSubmittedState{s},B_{t_s}),
    &\text{commit}.
 \end{cases}
 \label{app:eq:operation-action}
\end{equation}
Here $\mathsf{query}$ denotes a probe, $J_s$ is its requested readout set
($S$ in the probe rule), with
$J_s\subseteq\MaskSet{\OpSubmittedState{s}}$.
At the current committed history $(G,x_G)$, a probe state satisfies
\[
 (\OpSubmittedState{s})_i=x_i\quad(i\in G),
 \qquad
 (\OpSubmittedState{s})_i\in\Vocab\cup\{\MASK\}\quad(i\notin G).
\]
The decision rule in \cref{app:eq:operation-action} selects these
uncommitted entries and $J_s$ from the visible history and fixed seed;
any token assignments there are temporary and may all be omitted.
At a commit, the state instead has all uncommitted positions masked:
\[
 \OpSubmittedState{s}=y_{t_s},
 \qquad
 (y_{t_s})_i=
 \begin{cases}
  x_i,&i\in G,\\
  \MASK,&i\notin G,
 \end{cases}
\]
as in \cref{app:eq:commit-state}, with $G=G_{t_s-1}$.

% Please keep the role table with its introductory sentence.
\begin{samepage}
At the current committed history $(G,x_G)$,
the vertex roles are:
\begin{center}
\begin{tabular}{@{}lll@{}}
\toprule
Role & Probe & Commit $t$\\
\midrule
Temporarily assigned positions
 & $\ObservedSet{\OpSubmittedState{s}}\setminus G$ & $\varnothing$\\
Requested readouts & $J_s$ & $B_t$\\
Newly committed positions & $\varnothing$ & $B_t$\\
\bottomrule
\end{tabular}
\end{center}
\end{samepage}
If $z_s$ is the vector drawn at a commit, the ordinary observation and
the visible-history update are
\begin{equation}
 \OpVisibleObservation{s}:=
 \begin{cases}
  (\OracleRow{\OpSubmittedState{s}}{j})_{j\in J_s},
    &\text{query},\\
  ((\OracleRow{\OpSubmittedState{s}}{j})_{j\in B_{t_s}},z_s),
    &\text{commit},
 \end{cases}
 \label{app:eq:operation-visible-observation}
\end{equation}
\begin{equation}
 \OpVisibleHistory{(s+1)}
 :=(\OpVisibleHistory{s},A_s,\OpVisibleObservation{s}).
 \label{app:eq:operation-visible-history}
\end{equation}
The decision-rule internal state is a deterministic function of this
history and $w$. This full-history representation does not require the
implementation to store the entire transcript.
\end{definition}

\begin{definition}[\textcolor{red}{Output law and \HellingerCase\ risk}]
\label{app:def:output-risk}
\textcolor{red}{For an admissible algorithm from
\cref{app:def:admissible-algorithm}, let $\SamplerOutput$ be its output.}
Conditional on $\ControllerSeed=w$, let
$\SamplerOutputLaw{w}{\FrozenOracle}$ be its law after integrating all
product-commit draws.
Define the seed-averaged total-variation risk by
\begin{equation}
 \SeedRisk(\mathcal A;\TargetLaw,\FrozenOracle)
 :=\mathbb E_{\ControllerSeed}
 \dTV\bigl(\TargetLaw,
            \SamplerOutputLaw{\ControllerSeed}{\FrozenOracle}\bigr),
 \label{app:eq:seed-risk}
\end{equation}
and the oracle-robust risk by
\begin{equation}
 \RobustRisk(\mathcal A;\TargetLaw)
 :=\sup_{\FrozenOracle\in\OracleClass(\TargetLaw;\OracleRadius)}
 \SeedRisk(\mathcal A;\TargetLaw,\FrozenOracle).
 \label{app:eq:robust-risk}
\end{equation}
The expectation in \eqref{app:eq:seed-risk} is the expectation of the
conditional-output TV, rather than the TV after first mixing over
$\ControllerSeed$.
In \KLCase, the same conditional output law is evaluated by the
seed-averaged forward-KL risk of \cref{app:def:kl-oracle-risk}.
\end{definition}

\begin{definition}[\textcolor{red}{Resources and all-path budgets}]
\label{app:def:resources}
\textcolor{red}{These definitions apply to both accuracy cases.
In the counts below, $s$ ranges over the positive-integer indices of
submissions actually performed, as in \cref{app:eq:operation-action};
$t$ indexes commits only.}
Let $s_{\mathrm{pre}}$ be the index of the one distinguished preprocessing
all-mask readout, with $s_{\mathrm{pre}}=0$ if it is not used. Define
\[
\begin{aligned}
 \CounterfactualQueries
 &:=\bigl|\{s:s\ne s_{\mathrm{pre}},\
                 \text{submission }s\text{ is a probe}\}\bigr|,\\
 \CommitRounds
 &:=\bigl|\{s:\text{submission }s\text{ performs a nonempty product commit}\}\bigr|.
\end{aligned}
\]
Each probe other than that preprocessing readout contributes one to
$\CounterfactualQueries$; each commit
contributes one to $\CommitRounds$ and zero to $\CounterfactualQueries$,
regardless of its batch size. A probe at the commit state
\eqref{app:eq:commit-state} still counts toward $\CounterfactualQueries$:
the operation performed, rather than the masked state alone, determines
its counter. Every other all-mask probe also counts toward
$\CounterfactualQueries$. Reusing returned rows in several tests adds
no submission, whereas submitting the same state again does.
The actual total number of submissions is therefore
$\CounterfactualQueries+\CommitRounds+\mathbf 1\{s_{\mathrm{pre}}>0\}$.

The one distinguished all-mask readout is kept separate through the
accounting flag $\MarginalQueryCharge\in\{0,1\}$: set it to one when
that readout is charged, and to zero when it is omitted from the reported
count (or is not used). The reported \emph{noncommit} count is
\begin{equation}
 \ReportedQueries
 :=\CounterfactualQueries+\MarginalQueryCharge.
 \label{app:eq:reported-query-count}
\end{equation}
Thus $\ReportedQueries$ excludes commits. When the preprocessing readout
is used and charged, the full count is
$\ReportedQueries+\CommitRounds=1+\CounterfactualQueries+\CommitRounds$,
as in the main construction. Setting $\MarginalQueryCharge=0$ omits only
that preprocessing charge; it does not make other probes free.

\emph{Sequential depth and bandwidth.}
$\OracleDepth$ counts sequential oracle stages. A stage may contain
multiple submissions whose inputs are fixed before receiving that
stage's replies; later inputs that depend on those replies require
another stage. Parallel submissions therefore share depth but each
retains its own submission charge. Returned-row counts, scalar bandwidth,
and local computation are not included in these interaction counts;
\cref{app:lem:upper-resources} separately accounts for the construction's rows.

\smallskip
\noindent\textbf{Deterministic all-path budgets.}
The budgets $\QueryBudget,\RoundBudget$ satisfy
$\CounterfactualQueries\le\QueryBudget$ and $\CommitRounds\le\RoundBudget$
on every instance, seed, transcript, and realized path.  A lower bound on
these budgets does not assert the same lower bound on every realized counter.
\end{definition}

\begin{example*}[Submission counting]
One preprocessing all-mask readout, two probes of the same state, and
one commit of all $N$ positions use
$\CounterfactualQueries=2$ and $\CommitRounds=1$.
Charging preprocessing gives $\MarginalQueryCharge=1$,
$\ReportedQueries=3$, and $\ReportedQueries+\CommitRounds=4$ evaluations.
Returning several rows per probe or using each row in several tests
does not change these counts.
\end{example*}

\paragraph{Quantifier order.}
The public class data and the algorithm
are fixed first.  A hidden target law satisfying
Assumptions~\ref{app:ass:forest-structure} and
\ref{app:ass:response-regularity}, together with
Assumption~\ref{app:ass:rate-regime} where asymptotic rates are claimed, is
then chosen. In \HellingerCase, a single frozen oracle satisfying
Assumption~\ref{app:ass:oracle-accuracy} is chosen after the target law; the algorithm
is then executed and the robust risk in \eqref{app:eq:robust-risk} is evaluated.
\KLCaseTitle\ uses the same order with the oracle class and output risk
of \cref{app:def:kl-oracle-risk}.
